%%%%%%%%%%%%%%%%
%%% lev 5 %%%
%%%%%%%%%%%%%%%%

\Chapter{5}

\Verse{1}
{
	\hP{וְנֶפֶשׁ כִּי תֶחֱטָא וְשָׁמְעָה קוֹל אָלָה וְהוּא עֵד אוֹ רָאָה אוֹ יָדָע אִם לוֹא יַגִּיד וְנָשָׂא עֲוֺנוֹ׃}
}
{
	\eP{
		5 “And a person who will sin in that he has heard a
		pronouncement of an oath, and he was a witness—whether he
		saw or he knew: if he will not tell, then he shall bear
		his crime.
	}
}
{
%	\footnote{
%		bear his crime. This expression means that the person is not punished by
%		a court. It is between the individual and God. After all, in this case,
%		it may be that no one but the person himself is even aware that he did
%		the sin. The reverse expression occurs in the great formula of God’s
%		compassion in Exod 34:7, where it says that God bears crime (ns’ ‘wn)
%		for humans. That very emphasis on the usual divine mercy makes cases
%		like the current one so frightening: here it is the human who must bear
%		the crime himself. This expression will be used in a number of cases in
%		Leviticus and Numbers. By saying that a case is beyond the jurisdiction
%		of any court or ruler, the Torah recognizes that some things are beyond
%		humans’ ability and humans’ right to judge or punish.
%	}
}

\Verse{2}
{
	\hP{אוֹ נֶפֶשׁ אֲשֶׁר תִּגַּע בְּכל דָּבָר טָמֵא אוֹ בְנִבְלַת חַיָּה טְמֵאָה אוֹ בְּנִבְלַת בְּהֵמָה טְמֵאָה אוֹ בְּנִבְלַת שֶׁרֶץ טָמֵא וְנֶעְלַם מִמֶּנּוּ וְהוּא טָמֵא וְאָשֵׁם׃}
}
{
	\eP{
		“Or a person who will touch any impure thing or the
		carcass of an impure wild animal or the carcass of an
		impure domestic animal or the carcass of an impure
		swarming creature, and it was hidden from him, so he had
		become impure and had become guilty;
	}
}
{
}

\Verse{3}
{
	\hP{אוֹ כִי יִגַּע בְּטֻמְאַת אָדָם לְכֹל טֻמְאָתוֹ אֲשֶׁר יִטְמָא בָּהּ וְנֶעְלַם מִמֶּנּוּ וְהוּא יָדַע וְאָשֵׁם׃}
}
{
	\eP{
		or when he will touch a human’s impurity—for any impurity
		of his through which he will become impure—and it was
		hidden from him, and then he had come to know and had
		become guilty;
	}
}
{
}

\Verse{4}
{
	\hP{אוֹ נֶפֶשׁ כִּי תִשָּׁבַע לְבַטֵּא בִשְׂפָתַיִם לְהָרַע אוֹ לְהֵיטִיב לְכֹל אֲשֶׁר יְבַטֵּא הָאָדָם בִּשְׁבֻעָה וְנֶעְלַם מִמֶּנּוּ וְהוּא יָדַע וְאָשֵׁם לְאַחַת מֵאֵלֶּה׃}
}
{
	\eP{
		or a person who will swear so as to let out of his lips to
		do bad or to do good—for anything that a human would let
		out in an oath—and it was hidden from him, and then he had
		come to know and had become guilty by one of these:
	}
}
{
}

\Verse{5}
{
	\hP{וְהָיָה כִי יֶאְשַׁם לְאַחַת מֵאֵלֶּה וְהִתְוַדָּה אֲשֶׁר חָטָא עָלֶיהָ׃}
}
{
	\eP{
		And it will be that when he becomes guilty by one of these
		he shall confess that he has sinned over it,
	}
}
{
}

\Verse{6}
{
	\hP{וְהֵבִיא אֶת אֲשָׁמוֹ לַיהֹוָה עַל חַטָּאתוֹ אֲשֶׁר חָטָא נְקֵבָה מִן הַצֹּאן כִּשְׂבָּה אוֹ שְׂעִירַת עִזִּים לְחַטָּאת וְכִפֶּר עָלָיו הַכֹּהֵן מֵחַטָּאתוֹ׃}
}
{
	\eP{
		and he shall bring his guilt offering to YHWH over his sin
		that he has committed: a female from the flock, a sheep or
		a goat, as a sin offering, and the priest shall make
		atonement over him from his sin.
	}
}
{
}

\Verse{7}
{
	\hP{וְאִם לֹא תַגִּיעַ יָדוֹ דֵּי שֶׂה וְהֵבִיא אֶת אֲשָׁמוֹ אֲשֶׁר חָטָא שְׁתֵּי תֹרִים אוֹ שְׁנֵי בְנֵי יוֹנָה לַיהֹוָה אֶחָד לְחַטָּאת וְאֶחָד לְעֹלָה׃}
}
{
	\eP{
		“And if his hand will not attain enough for a sheep, then
		he shall bring as his guilt offering for having sinned two
		turtledoves or two pigeons to YHWH: one for a sin offering
		and one for a burnt offering.
	}
}
{
%	\footnote{
%		if his hand will not attain enough. Meaning: if he will not be able to
%		afford it. A person is not to be prevented from getting atonement
%		because of lack of money.
%	}
}

\Verse{8}
{
	\hP{וְהֵבִיא אֹתָם אֶל הַכֹּהֵן וְהִקְרִיב אֶת אֲשֶׁר לַחַטָּאת רִאשׁוֹנָה וּמָלַק אֶת רֹאשׁוֹ מִמּוּל ערְפּוֹ וְלֹא יַבְדִּיל׃}
}
{
	\eP{
		And he shall bring them to the priest, and he shall bring
		forward the one that is for the sin offering first and
		will wring but not separate its head at the point opposite
		its neck.
	}
}
{
}

\Verse{9}
{
	\hP{וְהִזָּה מִדַּם הַחַטָּאת עַל קִיר הַמִּזְבֵּחַ וְהַנִּשְׁאָר בַּדָּם יִמָּצֵה אֶל יְסוֹד הַמִּזְבֵּחַ חַטָּאת הוּא׃}
}
{
	\eP{
		And he shall sprinkle some of the blood of the sin
		offering on the wall of the altar, and what remains of the
		blood shall be drained at the base of the altar. It is a
		sin offering.
	}
}
{
}

\Verse{10}
{
	\hP{וְאֶת הַשֵּׁנִי יַעֲשֶׂה עֹלָה כַּמִּשְׁפָּט וְכִפֶּר עָלָיו הַכֹּהֵן מֵחַטָּאתוֹ אֲשֶׁר חָטָא וְנִסְלַח לוֹ׃}
}
{
	\eP{
		And he shall make the second one a burnt offering
		according to the required manner. And the priest shall
		make atonement over him from his sin that he committed,
		and it shall be forgiven for him.
	}
}
{
}

\Verse{11}
{
	\hP{וְאִם לֹא תַשִּׂיג יָדוֹ לִשְׁתֵּי תֹרִים אוֹ לִשְׁנֵי בְנֵי יוֹנָה וְהֵבִיא אֶת קרְבָּנוֹ אֲשֶׁר חָטָא עֲשִׂירִת הָאֵפָה סֹלֶת לְחַטָּאת לֹא יָשִׂים עָלֶיהָ שֶׁמֶן וְלֹא יִתֵּן עָלֶיהָ לְבֹנָה כִּי חַטָּאת הִוא׃}
}
{
	\eP{
		“And if his hand cannot attain enough for two turtledoves
		or two pigeons, then he shall bring as his offering for
		having sinned a tenth of an ephah of fine flour for a sin
		offering. He shall not set oil on it, and he shall not put
		frankincense on it, because it is a sin offering.
	}
}
{
}

\Verse{12}
{
	\hP{וֶהֱבִיאָהּ אֶל הַכֹּהֵן וְקָמַץ הַכֹּהֵן מִמֶּנָּה מְלוֹא קֻמְצוֹ אֶת אַזְכָּרָתָהּ וְהִקְטִיר הַמִּזְבֵּחָה עַל אִשֵּׁי יְהֹוָה חַטָּאת הִוא׃}
}
{
	\eP{
		And he shall bring it to the priest, and the priest shall
		take the fill of his fist from it, a representative
		portion of it, and he shall burn it to smoke at the altar
		with YHWH’s offerings by fire. It is a sin offering.
	}
}
{
}

\Verse{13}
{
	\hP{וְכִפֶּר עָלָיו הַכֹּהֵן עַל חַטָּאתוֹ אֲשֶׁר חָטָא מֵאַחַת מֵאֵלֶּה וְנִסְלַח לוֹ וְהָיְתָה לַכֹּהֵן כַּמִּנְחָה׃}
}
{
	\eP{
		And the priest shall make atonement over him, over his sin
		that he committed, from one of these, and it shall be
		forgiven for him. And it shall be the priest’s, like the
		grain offering.”
	}
}
{
}

\Verse{14}
{
	\hP{וַיְדַבֵּר יְהֹוָה אֶל מֹשֶׁה לֵּאמֹר׃}
}
{
	\eP{And YHWH spoke to Moses, saying,}
}
{
}

\Verse{15}
{
	\hP{נֶפֶשׁ כִּי תִמְעֹל מַעַל וְחָטְאָה בִּשְׁגָגָה מִקּדְשֵׁי יְהֹוָה וְהֵבִיא אֶת אֲשָׁמוֹ לַיהֹוָה אַיִל תָּמִים מִן הַצֹּאן בְּעֶרְכְּךָ כֶּסֶף שְׁקָלִים בְּשֶׁקֶל הַקֹּדֶשׁ לְאָשָׁם׃}
}
{
	\eP{
		“A person who will make a breach and sin by mistake among
		YHWH’s holy things: he shall bring his guilt offering to
		YHWH, an unblemished ram from the flock by your evaluation
		of silver by shekels, by the shekel of the Holy, as a
		guilt offering.
	}
}
{
%	\footnote{
%		make a breach. The term, Hebrew ma‘al, refers to trespassing into a zone
%		that is forbidden. Here it refers to a layperson’s mistakenly taking
%		some of the holy objects, something that is no longer available, which
%		would in most cases mean food that was meant for the priests. (Lev
%		22:14–16 refers explicitly to a layperson mistakenly eating the holy
%		things, and the penalty is the same: add a fifth.) In Joshua 7:1, this
%		term, “breach,” refers to Achan’s taking some of the spoils of Jericho,
%		which the Israelites had been forbidden to take under the law of erem
%		(complete destruction in wars commanded by God; see, e.g., Deut 13:18).
%		In Num 5:6, it refers to unlawful acquisition of another person’s
%		property. In Num 5:12, a woman’s adultery is referred to as a ma‘al
%		against her husband. In Lev 26:40, it refers generally to a breach of
%		faith toward God by the community.
%	}
}

\Verse{16}
{
	\hP{וְאֵת אֲשֶׁר חָטָא מִן הַקֹּדֶשׁ יְשַׁלֵּם וְאֶת חֲמִישִׁתוֹ יוֹסֵף עָלָיו וְנָתַן אֹתוֹ לַכֹּהֵן וְהַכֹּהֵן יְכַפֵּר עָלָיו בְּאֵיל הָאָשָׁם וְנִסְלַח לוֹ׃}
}
{
	\eP{
		And for what he sinned from the Holy he shall pay and add
		to it a fifth of it and give it to the priest, and the
		priest shall make atonement over him with the ram of the
		guilt offering, and it shall be forgiven for him.
	}
}
{
}

\Verse{17}
{
	\hP{וְאִם נֶפֶשׁ כִּי תֶחֱטָא וְעָשְׂתָה אַחַת מִכּל מִצְוֺת יְהֹוָה אֲשֶׁר לֹא תֵעָשֶׂינָה וְלֹא יָדַע וְאָשֵׁם וְנָשָׂא עֲוֺנוֹ׃}
}
{
	\eP{
		“And if a person who will sin and commit one of any of
		YHWH’s commandments that are not to be done and did not
		know and became guilty, then he shall bear his crime.
	}
}
{
}

\Verse{18}
{
	\hP{וְהֵבִיא אַיִל תָּמִים מִן הַצֹּאן בְּעֶרְכְּךָ לְאָשָׁם אֶל הַכֹּהֵן וְכִפֶּר עָלָיו הַכֹּהֵן עַל שִׁגְגָתוֹ אֲשֶׁר שָׁגָג וְהוּא לֹא יָדַע וְנִסְלַח לוֹ׃}
}
{
	\eP{
		And he shall bring an unblemished ram from the flock by
		your evaluation as a guilt offering to the priest, and the
		priest shall make atonement over him, over his mistake
		that he made and had not known, and it shall be forgiven
		for him.
	}
}
{
}

\Verse{19}
{
	\hP{אָשָׁם הוּא אָשֹׁם אָשַׁם לַיהֹוָה׃}
}
{
	\eP{It is a guilt offering. He is guilty to YHWH.”}
}
{
}

\Verse{20}
{
	\hP{וַיְדַבֵּר יְהֹוָה אֶל מֹשֶׁה לֵּאמֹר׃}
}
{
	\eP{And YHWH spoke to Moses, saying,}
}
{
}

\Verse{21}
{
	\hP{נֶפֶשׁ כִּי תֶחֱטָא וּמָעֲלָה מַעַל בַּיהֹוָה וְכִחֵשׁ בַּעֲמִיתוֹ בְּפִקָּדוֹן אוֹ בִתְשׂוּמֶת יָד אוֹ בְגָזֵל אוֹ עָשַׁק אֶת עֲמִיתוֹ׃}
}
{
	\eP{
		“A person who will sin and make a breach against YHWH and
		tell a lie against his fellow in a matter of a deposit or
		something set in hand or by robbery, or has exploited his
		fellow
	}
}
{
}

\Verse{22}
{
	\hP{אוֹ מָצָא אֲבֵדָה וְכִחֶשׁ בָּהּ וְנִשְׁבַּע עַל שָׁקֶר עַל אַחַת מִכֹּל אֲשֶׁר יַעֲשֶׂה הָאָדָם לַחֲטֹא בָהֵנָּה׃}
}
{
	\eP{
		or has found something that was lost and has lied about it
		or sworn to a falsehood—for one of any of these that a
		human will do to sin:
	}
}
{
}

\Verse{23}
{
	\hP{וְהָיָה כִּי יֶחֱטָא וְאָשֵׁם וְהֵשִׁיב אֶת הַגְּזֵלָה אֲשֶׁר גָּזָל אוֹ אֶת הָעֹשֶׁק אֲשֶׁר עָשָׁק אוֹ אֶת הַפִּקָּדוֹן אֲשֶׁר הפְקַד אִתּוֹ אוֹ אֶת הָאֲבֵדָה אֲשֶׁר מָצָא׃}
}
{
	\eP{
		it will be, when he sins and is guilty, that he shall
		bring back the thing that he robbed or the thing that he
		coerced or the thing that was deposited with him or the
		lost thing that he found
	}
}
{
}

\Verse{24}
{
	\hP{אוֹ מִכֹּל אֲשֶׁר יִשָּׁבַע עָלָיו לַשֶּׁקֶר וְשִׁלַּם אֹתוֹ בְּרֹאשׁוֹ וַחֲמִשִׁתָיו יֹסֵף עָלָיו לַאֲשֶׁר הוּא לוֹ יִתְּנֶנּוּ בְּיוֹם אַשְׁמָתוֹ׃}
}
{
	\eP{
		or anything about which he swore as a falsehood, and he
		shall pay it in its worth and add to it a fifth of it. He
		shall give it to the one whose it is on the day of his
		being guilty,
	}
}
{
}

\Verse{25}
{
	\hP{וְאֶת אֲשָׁמוֹ יָבִיא לַיהֹוָה אַיִל תָּמִים מִן הַצֹּאן בְּעֶרְכְּךָ לְאָשָׁם אֶל הַכֹּהֵן׃}
}
{
	\eP{
		and he shall bring his guilt offering to YHWH: an
		unblemished ram from the flock by your evaluation as a
		guilt offering, to the priest.
	}
}
{
}

\Verse{26}
{
	\hP{וְכִפֶּר עָלָיו הַכֹּהֵן לִפְנֵי יְהֹוָה וְנִסְלַח לוֹ עַל אַחַת מִכֹּל אֲשֶׁר יַעֲשֶׂה לְאַשְׁמָה בָהּ׃}
}
{
	\eP{
		And the priest shall atone over him in front of YHWH, and
		it will be forgiven for him, for any one out of all that
		he will do to become guilty through it.”
	}
}
{
}
